\documentclass[sigconf,nonacm]{acmart}

\usepackage{booktabs}
\usepackage{listings}
\usepackage{xcolor}
\usepackage{graphicx}
\usepackage{subcaption}

\lstset{
  basicstyle=\ttfamily\small,
  keywordstyle=\color{blue},
  commentstyle=\color{gray},
  stringstyle=\color{red},
  breaklines=true,
  frame=single,
}

\title{Deep Semantic Test Intelligence: Extracting Method Intent Signals for Automated Test Generation}

\author{DocSpec Research Team}
\affiliation{%
  \institution{DocSpec Project}
}
\email{research@docspec.io}

\begin{abstract}
We present Deep Semantic Test Intelligence (DSTI), a novel approach to
extracting semantic intent signals from source code across 13 orthogonal
channels, including naming semantics, guard clauses, branching patterns,
data flow, error handling, and validation annotations. DSTI computes an
Intent Signal Density (ISD) score that quantifies method complexity and
guides automated test generation. We implement DSTI in DocSpec, a
universal documentation specification tool, and evaluate it against 835
real-world bugs from the Defects4J benchmark across 17 Java projects.
Our results show that DSTI-generated test stubs achieve [X]\% mutation
coverage improvement over baseline approaches, with an average analysis
time of [Y] seconds per method. A developer study with 30 participants
demonstrates that DSTI signals reduce time-to-first-test by [Z]\% and
improve test relevance ratings by [W] points on a 5-point scale.
\end{abstract}

\keywords{test generation, semantic analysis, documentation, intent signals}

\begin{document}

\maketitle

\section{Introduction}

Software testing remains one of the most labor-intensive activities in
software development. Despite advances in automated test generation tools
such as Randoop~\cite{pacheco2007randoop} and EvoSuite~\cite{fraser2011evosuite},
generated tests often lack semantic relevance---they exercise code paths
but fail to capture the developer's intended behavior.

We observe that source code contains rich \emph{intent signals} beyond
its executable semantics. Method names encode purpose (``validate,''
``calculate,'' ``transform''), guard clauses express preconditions,
branching patterns reveal decision logic, and validation annotations
declare constraints. These signals, while implicit, form a comprehensive
picture of what a method \emph{should} do.

This paper makes the following contributions:
\begin{enumerate}
  \item \textbf{DSTI Framework:} A systematic taxonomy of 13 intent signal
    channels extracted from source code via static analysis.
  \item \textbf{ISD Score:} A weighted composite metric that quantifies
    intent density, enabling prioritization of test generation effort.
  \item \textbf{Property DSL:} A domain-specific language for expressing
    method invariants, enabling property-based test generation.
  \item \textbf{Evaluation:} Empirical evaluation on Defects4J and a
    controlled developer study.
\end{enumerate}

\section{Background and Related Work}

\subsection{Automated Test Generation}
% EvoSuite, Randoop, Pex, IntelliTest
% Limitations: semantic relevance, oracle problem

\subsection{Specification Mining}
% Daikon, JML, contracts
% Difference: DSTI extracts from code structure, not runtime traces

\subsection{Documentation Generation}
% Javadoc, Swagger/OpenAPI, TypeDoc
% Limitation: API surface only, no behavioral intent

\section{Deep Semantic Test Intelligence}

\subsection{Intent Signal Channels}

DSTI extracts signals across 13 orthogonal channels:

\begin{enumerate}
  \item \textbf{Name Semantics} --- verb/object decomposition of method names
  \item \textbf{Guard Clauses} --- precondition checks at method entry
  \item \textbf{Branches} --- conditional logic patterns
  \item \textbf{Data Flow} --- field reads and writes
  \item \textbf{Return Type} --- complexity of the return type
  \item \textbf{Loops} --- iteration patterns (streams, enhanced for)
  \item \textbf{Error Handling} --- try/catch blocks and exception types
  \item \textbf{Constants} --- magic numbers and string literals
  \item \textbf{Null Checks} --- null-safety patterns
  \item \textbf{Assertions} --- programmatic assertions
  \item \textbf{Logging} --- log statements indicating intent
  \item \textbf{Dependencies} --- external service calls
  \item \textbf{Validation Annotations} --- Bean Validation, etc.
\end{enumerate}

\subsection{Intent Signal Density (ISD)}

The ISD score combines channel signals into a single metric:

\begin{equation}
  \text{ISD} = \sum_{i=1}^{13} w_i \cdot s_i(m)
\end{equation}

where $w_i$ is the channel weight and $s_i(m)$ is the normalized signal
strength for method $m$ on channel $i$. Weights are empirically tuned
(Table~\ref{tab:weights}).

\subsection{Property DSL}

The Property DSL enables developers to express invariants declaratively:

\begin{lstlisting}[language=Java]
@DocInvariant(rules = {
    "balance >= 0",
    "transactions SIZE > 0",
    "currency IN [USD, EUR, GBP]",
    "amount BETWEEN 0.01 AND 1000000"
})
\end{lstlisting}

\section{Implementation}

DocSpec implements DSTI as a Java annotation processor that runs during
compilation. The processor performs a single-pass AST traversal,
collecting intent signals from each method, computing ISD scores, and
generating both documentation (docspec.json) and test stubs.

\section{Evaluation}

\subsection{Defects4J Benchmark}
% Setup, methodology, results

\subsection{Developer Study}
% Protocol, results

\section{Results}

% Placeholder for actual results

\section{Threats to Validity}

\subsection{Internal Validity}
% Learning effects mitigated by crossover design

\subsection{External Validity}
% Java-only in current evaluation

\subsection{Construct Validity}
% ISD weights are empirically tuned

\section{Conclusion}

% Summary and future work

\bibliographystyle{ACM-Reference-Format}
\bibliography{references}

\end{document}
